\section{Regular expressions and Kleene's theorem}

\subsection{Regular expressions as shorthand}
A regular expression is a compact description of a regular
language. Examples over $\Sigma=\set{0,1}$:
\[
\begin{aligned}
  0^*10^*
    &=\set{w:w\text{ contains exactly one }1},\\
  \Sigma^*1
    &=\set{w:w\text{ ends with }1},\\
  0\Sigma^*\cup\Sigma^*1
    &=\set{w:w\text{ starts with }0\text{ or ends with }1},\\
  (0\cup 1)^*001(0\cup 1)^*
    &=\set{w:w\text{ contains }001\text{ as a substring}},\\
  (00)^*
    &=\set{w:|w|\text{ is even and }w\in 0^*},\\
  (0\cup 1)(0\cup 1)
    &=\set{w:|w|=2},\\
  (0\cup\varepsilon)1^*
    &=\set{0^i1^k:i\in\set{0,1},\,k\geq 0},\\
  0(0\cup 1)^*0\cup 1(0\cup 1)^*1\cup 0\cup 1
    &=\set{w:w\text{ starts and ends with the same symbol}}.
\end{aligned}
\]
Useful identities include
\[
  A\cup\varnothing=A,\qquad
  A\concat\varnothing=\varnothing,\qquad
  A\concat\set{\eps}=A,\qquad
  A\cup\set{\eps}=A\text{ if and only if }\eps\in A.
\]
$R\cup\eps$ need not equal $R$ unless $\eps$ is already in $\lang{R}$.

\textbf{Operator precedence} (highest to lowest): star, then concatenation, then
union. Use parentheses to override.

\subsection{Formal definition}
\begin{definition}[Regular expression]
$R$ is a regular expression over $\Sigma$ if it has one of the forms
\begin{enumerate}
  \item $R=a$ for some $a\in\Sigma$;
  \item $R=\eps$;
  \item $R=\varnothing$;
  \item $R=(R_1\cup R_2)$;
  \item $R=(R_1R_2)$;
  \item $R=(R_1^*)$;
\end{enumerate}
where $R_1$ and $R_2$ are regular expressions. The language $\lang{R}$ is
defined inductively:
\[
\begin{aligned}
  \lang{a}&=\set{a},&
  \lang{\eps}&=\set{\eps},&
  \lang{\varnothing}&=\varnothing,\\
  \lang{R_1\cup R_2}&=\lang{R_1}\cup\lang{R_2},&
  \lang{R_1R_2}&=\lang{R_1}\concat\lang{R_2},&
  \lang{R_1^*}&=\lang{R_1}^*.
\end{aligned}
\]
\end{definition}
Write $R^+$ for $RR^*$, which excludes $\eps$ unless $\eps\in\lang{R}$. Then
$R^*\equiv R^+\cup\set{\eps}$.

\subsection{Kleene's theorem}
\begin{theorem}[Kleene]
A language is regular if and only if some regular expression describes it.
\end{theorem}
The two directions are the next two lemmas.

\begin{lemma}[Expression $\Rightarrow$ automaton]
If a language is described by a regular expression, then it is regular.
\end{lemma}

\begin{proof}
Let $R$ be a regular expression. Induct on the structure of $R$.
\begin{enumerate}
  \item $R=a$: a two-state NFA with a single $a$-edge from start to accept.
  \item $R=\eps$: a single accepting start state and no edges (or an
        $\eps$-edge to an accept state).
  \item $R=\varnothing$: a start state that is not accepting, and no edges.
  \item $R=R_1\cup R_2$, $R=R_1R_2$, $R=R_1^*$: apply the NFA constructions of
        Chapter~\ref{ch:nfa} to the machines given by the inductive hypothesis.
\end{enumerate}
\end{proof}

\begin{exercise}
Build, by the inductive method, an NFA for $(ab\cup a)^*$ and compare it with
an NFA for $(ab^*\cup a^*)$. They are not the same language.
\end{exercise}

\subsection{Generalized NFAs}
\begin{lemma}[Automaton $\Rightarrow$ expression]
If a language is regular, then it is described by a regular expression.
\end{lemma}

\begin{proof}[Idea]
Convert a DFA to a \emph{generalized NFA} (GNFA), then eliminate states until
a single regular-expression label remains.

A GNFA has regular expressions (not just symbols) as edge labels, and satisfies:
\begin{enumerate}
  \item a unique start state with edges to every other state, and no incoming
        edges;
  \item a unique accept state with edges from every other state, and no
        outgoing edges;
  \item for every pair of non-start, non-accept states (including a state and
        itself), edges in both directions.
\end{enumerate}
Missing edges are labelled $\varnothing$.

\textbf{DFA to GNFA.} Add a new start state with an $\eps$-edge to the old
start. Add a new unique accept state with $\eps$-edges from every old accept
state. Combine multiple labels $R,S$ on the same pair of states into
$R\cup S$. Add $\varnothing$ labels where no edge existed.

\begin{definition}[GNFA]
A GNFA is a $5$-tuple $(Q,\Sigma,\transition,q_{\mathrm{start}},q_{\mathrm{accept}})$
where
\[
  \transition:
    \bigl(Q\setminus\set{q_{\mathrm{accept}}}\bigr)
    \times
    \bigl(Q\setminus\set{q_{\mathrm{start}}}\bigr)
    \to \mathcal{R}
\]
and $\mathcal{R}$ is the set of regular expressions over $\Sigma$. It accepts $w$
if $w$ can be broken into pieces matching the labels of some path from
$q_{\mathrm{start}}$ to $q_{\mathrm{accept}}$.
\end{definition}

\begin{center}
\begin{tikzpicture}[automaton,node distance=28mm]
  \node[state] (qi) {$q_i$};
  \node[state,right=of qi] (rip) {$q_{\mathrm{rip}}$};
  \node[state,right=of rip] (qj) {$q_j$};
  \path[->]
    (qi) edge[bend left=18] node[above]{$R_1$} (rip)
    (rip) edge[loop above] node{$R_2$} ()
    (rip) edge[bend left=18] node[above]{$R_3$} (qj)
    (qi) edge[bend right=28] node[below]{$R_4$} (qj);
\end{tikzpicture}
\end{center}
After removing $q_{\mathrm{rip}}$, the direct label $q_i\to q_j$ becomes
$R_4\cup R_1(R_2)^*R_3$.

\textbf{State elimination.} While more than two states remain, pick a state
$q_{\mathrm{rip}}\notin\set{q_{\mathrm{start}},q_{\mathrm{accept}}}$ and remove
it. For every remaining $q_i,q_j$ (not the accept as source, not the start as
target), replace the $q_i\to q_j$ label $R_4$ by
\[
  R_4\cup\bigl(R_1(R_2)^*R_3\bigr),
\]
where $R_1$ labels $q_i\to q_{\mathrm{rip}}$, $R_2$ labels the loop at
$q_{\mathrm{rip}}$, $R_3$ labels $q_{\mathrm{rip}}\to q_j$, and $R_4$ is the old
direct label $q_i\to q_j$.

When two states remain, the label of the unique edge
$q_{\mathrm{start}}\to q_{\mathrm{accept}}$ is the desired expression.

Correctness is by induction on the number $k$ of states of the GNFA. The base
$k=2$ is immediate. The inductive step: any path
that used $q_{\mathrm{rip}}$ is accounted for by $R_1(R_2)^*R_3$, and paths
that never used it are accounted for by $R_4$.
\end{proof}

\begin{exercise}
Convert the odd-parity DFA and the substring-$001$ DFA to regular
expressions by state elimination.
\end{exercise}
